\newpage
\appendix
\begin{center}
	{\zihao{4}\bfseries 附\quad 录}
\end{center}

\section*{附录 A\quad 程序与数据文件说明}
\addcontentsline{toc}{section}{附录 A\quad 程序与数据文件说明}

本文的全部程序、数据与结果文件按“公共层—问题层—论文层”组织。其中公共层实现附录 2/3 的物理与通信公式，问题一至问题四共用，保证口径唯一；问题层按问题分文件组织，每问有独立入口脚本，可单独运行。

\begin{table}[htbp]
	\caption{程序与数据文件清单}
	\centering
	\zihao{-5}
	\begin{tabularx}{0.94\textwidth}{lX}
		\toprule
		文件 / 目录 & 说明 \\
		\midrule
		\texttt{code/plotstyle.py} & 全局绘图样式（SimSun 中文字体、统一调色板、DPI$\ge$300）与数据加载 \\
		\texttt{code/verify\_q2\_energy.py} & 问题二能耗与返航 SOC 的\textbf{独立复算}（逐段） \\
		\texttt{code/audit\_q2\_plan.py} & 问题二方案的独立可行性审计（能量 / 时间 / 资源 / 时限） \\
		\texttt{code/fig\_common.py} & 总体分析与公共模型章图（fig06--fig12） \\
		\texttt{code/fig\_q1.py} \texttt{fig\_q2.py} \texttt{fig\_q3.py} \texttt{fig\_q4.py} & 四问的图（fig13--fig35） \\
		\texttt{code/fig\_check.py} & 敏感性、独立复算审计与指标汇总图（fig36--fig38） \\
		\texttt{code/make\_flowcharts.py} & 生成四张求解流程图的 draw.io 源文件 \\
		\texttt{code/make\_tables.py} & 由 CSV 自动生成论文三线表，避免手工誊抄 \\
		\texttt{diagrams/*.drawio} & 技术路线图与四张流程图的\textbf{可编辑}源文件 \\
		\texttt{figures/} & 论文全部插图（PNG + PDF 矢量备份） \\
		\texttt{tables/} & 全部表格数据源（CSV） \\
		\texttt{data/} & 建模就绪数据、四问指标与参数、独立复算与审计结果 \\
		\texttt{results/} & 按赛题《结果提交模板》生成的 6 个交付 sheet \\
		\bottomrule
	\end{tabularx}
\end{table}

\section*{附录 B\quad 关键物理与通信口径汇总}
\addcontentsline{toc}{section}{附录 B\quad 关键物理与通信口径汇总}

为便于复核，此处集中列出本文全部关键口径。这些口径在四问中\textbf{只实现一次}，四问共用：

\begin{table}[htbp]
	\caption{关键物理与通信口径}
	\centering
	\zihao{-5}
	\begin{tabularx}{0.94\textwidth}{lX}
		\toprule
		项目 & 口径 \\
		\midrule
		巡航海拔 & 航段经过的 DEM 像元最高地面高程 $+50$\,m \\
		作业高度 & O01 取地面海拔；服务区取地面海拔 $+30$\,m \\
		载荷—航程 & $L_{g}(q)=L_{g0}-(L_{g0}-L_{gF})(q/Q_{g})^{3/2}$（指数 $3/2$） \\
		最大安全载荷 & 能量约束取等号，用单调性 $+$ Brent 法\textbf{数值反解}（无闭式解） \\
		水平巡航能耗 & 由航程反推 $E^{hor}=(d/L_{g}(q))E_{g}^{use}$（附件未给运输机巡航功率） \\
		爬升附加能耗 & $E^{up}=mgh^{+}/\eta_{up}$，$\eta_{up}=0.72$（附件列名为“爬升能耗效率”） \\
		下降能耗 & \textbf{不单独计}（附件下降能耗效率 $=0$） \\
		返航余量 & $E_{p}^{T}\le(1-\rho_{g})E_{g}^{use}$，$\rho_{g}=0.20$ \\
		充电模型 & 两阶段：SOC$<90\%$ 占 $T_{full}$ 的 65\%，$90\%\sim100\%$ 占 35\%（$s=0.9$ 处连续） \\
		FSPL & $32.45+20\log_{10}f(\mathrm{MHz})+20\log_{10}D(\mathbf{km})$ \\
		双向门限 & 取两个方向中的\textbf{较小值} \\
		通信三态 & 直连 $>$ 中继（接入段与回传段\textbf{同时}可用）$>$ 中断；\textbf{不允许多跳} \\
		多点载荷递减 & 第 $m$ 段携带“尚未投送的货”（飞到第 $k$ 站时仍带第 $k..M$ 站的货） \\
		交接时间 & 逐站累加 $\sum_{s}(t_{h0}+k_{s}t_{h1})$，\textbf{不}乘总箱数 \\
		\bottomrule
	\end{tabularx}
\end{table}

\section*{附录 C\quad 四问关键指标汇总}
\addcontentsline{toc}{section}{附录 C\quad 四问关键指标汇总}

\begin{table}[htbp]
	\caption{四问关键指标汇总}
	\centering
	\zihao{-5}
	\begin{tabular}{lll}
		\toprule
		问题 & 指标 & 数值 \\
		\midrule
		问题一 & 最优往返架次数 & 18（全部 C 型，逐区达解析下界） \\
		问题一 & 总运输能耗 & 75.07\,kWh \\
		问题一 & 累计作业时间 / 并行完工时间 & 9.37\,h / 1.27\,h \\
		问题一 & 最优性 & 架次数维度\textbf{可证最优}（15 区全部等于下界） \\
		\midrule
		问题二 & 运输架次数 & 25（A 型 8 $+$ B 型 10 $+$ C 型 7，8 架全用；理论下界 24，质量装载率 79\%） \\
		问题二 & 总运输能耗 & 77.30\,kWh \\
		问题二 & 全部任务完成时间 & 3.02\,h（10855.1\,s） \\
		问题二 & 期望送达准时率 & 100.0\%（80/80 箱全部按时送达；首批违规 0 箱、期望违规 0 箱） \\
		问题二 & 独立校验结论 & \textbf{违规 0 条}（物理类硬约束与两类时限约束全部通过） \\
		\midrule
		问题三 & 运输 / 中继架次数 & 25 / 20（其中 20 个运输架次需中继） \\
		问题三 & 需保障架次覆盖率 & 20 / 20 $=$ 100\%（均为单一悬停点全程覆盖） \\
		问题三 & 运输 / 中继 / 总能耗 & 77.30 / 4.83 / 82.14\,kWh（中继占 5.9\%） \\
		问题三 & 联合任务完成时间 & 3.15\,h（11323.3\,s） \\
		问题三 & 平均直连中断占比 & 28.4\% \\
		\midrule
		问题四 & 合法 2 / 3 组分区 & \textbf{均可零架次改动}（原子单元数 3：$K=2$ 与 $K=3$ 直接由原子单元合并得到） \\
		问题四 & 桥接架次 & 6 个（T019、T017、T018、T021、T024、T025） \\
		问题四 & 最小改动 & \textbf{0 个架次}（$K=2$ 与 $K=3$ 均为 0） \\
		问题四 & 资源总量（K$\,=\,$1/2/3） & 32 / 34 / 37（台$\cdot$组） \\
		问题四 & 组间不均衡度 & 0 / 1.8517 / 2.7035 \\
		问题四 & 最大组工作量（K$\,=\,$1/2/3） & 13.95 / 13.43 / 13.00\,h \\
		问题四 & K$\,=\,$1 资源缺口 & A、B、C 三型共享电池各缺 1 组（运输机、中继机与中继组件无缺口） \\
		\bottomrule
	\end{tabular}
\end{table}

\section*{附录 D\quad 可复现性说明}
\addcontentsline{toc}{section}{附录 D\quad 可复现性说明}

本文结果的可复现性由以下措施保证：

\begin{enumerate}[label=\arabic*.]
	\item \textbf{随机种子固定}：全部随机过程使用 SEED $=42$，同种子两次运行结果完全一致。
	\item \textbf{物理口径唯一}：附录 2/3 的公式只在公共层实现一次，四问共用，不存在口径分叉。
	\item \textbf{结果可追溯}：每个问题的指标、参数、图与表均有对应的输出文件；论文中的每个数值都能追溯到附件数据或程序输出。
	\item \textbf{独立复算}：问题二的能耗与 SOC 由第二套独立代码逐段重算验证，与求解器上报值一致（总能耗均为 77.3016\,kWh）；问题二、问题三的方案另由独立校验器逐条复核，违规 0 条。
	\item \textbf{图表自动生成}：全部插图由脚本生成（可复现），全部表格由 CSV 自动转 LaTeX（避免手工誊抄错误）。
\end{enumerate}

复现命令（Python 3.13，依赖 NumPy / SciPy / pandas / Matplotlib / seaborn / NetworkX / Plotly）：

\begin{lstlisting}[style=Python, basicstyle=\zihao{-5}\ttfamily, frame=single]
python code/verify_q2_energy.py   # 问题二能耗与 SOC 独立复算
python code/audit_q2_plan.py      # 问题二方案独立可行性审计
python code/make_flowcharts.py    # 生成技术路线图与四张流程图的 draw.io 源
python code/fig_common.py         # 总体分析与公共模型章图
python code/fig_q1.py             # 问题一图
python code/fig_q2.py             # 问题二图
python code/fig_q3.py             # 问题三图
python code/fig_q4.py             # 问题四图
python code/fig_check.py          # 敏感性、审计与指标汇总图
python code/make_tables.py        # 由 CSV 生成论文三线表
\end{lstlisting}
